\section{Training-Free GRPO Algorithm}

\label{sec:algorithm}

We now present the complete Training-Free Group Relative Policy Optimization algorithm.

\subsection{Overview}

Training-Free GRPO operates entirely in context space, using semantic advantages extracted via LLM introspection. The algorithm has three stages executed within each training iteration:

\begin{enumerate}
    \item \textbf{Stage 1: Trajectory Summarization} — Condense each rollout into a step-by-step natural language summary.
    \item \textbf{Stage 2: Group Advantage Extraction} — Compare trajectories within groups to identify strategic patterns and propose experience updates.
    \item \textbf{Stage 3: Batch Consolidation} — Merge and refine all proposed updates into final consolidated operations.
\end{enumerate}

\subsection{Algorithm Formulation}

\begin{algorithm}[H]
\caption{Training-Free Group Relative Policy Optimization}
\label{alg:tfgrpo}
\begin{algorithmic}[1]
\REQUIRE Training queries $\mathcal{Q} = \{q_1, \ldots, q_N\}$, ground truth answers $\mathcal{GT} = \{a_1, \ldots, a_N\}$, frozen foundation model $\pi_\theta$, group size $G$, number of epochs $T$
\ENSURE Experience library $\mathcal{E}$
\STATE Initialize $\mathcal{E} \leftarrow \emptyset$
\FOR{epoch $t = 1$ to $T$}
    \STATE $\text{AllOperations} \leftarrow []$
    \FOR{each query $q_i \in \mathcal{Q}$}
        \STATE \textcolor{blue}{// Stage 0: Generate group rollouts}
        \STATE $\text{Trajectories} \leftarrow []$
        \FOR{$g = 1$ to $G$}
            \STATE $\tau_g \leftarrow \pi_\theta(\cdot | q_i, \text{Format}(\mathcal{E}, q_i))$ \COMMENT{Context injection}
            \STATE $r_g \leftarrow \text{Reward}(\tau_g, a_i)$ \COMMENT{Evaluate correctness}
            \STATE $\text{Trajectories}.\text{append}((\tau_g, r_g))$
        \ENDFOR

        \IF{$\text{std}([r_1, \ldots, r_G]) = 0$}
            \STATE \textbf{continue} \COMMENT{Skip homogeneous groups}
        \ENDIF

        \STATE \textcolor{blue}{// Stage 1: Trajectory summarization}
        \STATE $\text{Summaries} \leftarrow []$
        \FOR{each $(\tau_g, r_g) \in \text{Trajectories}$}
            \STATE $s_g \leftarrow \text{LLM}_{\text{sum}}(\tau_g, r_g, a_i)$ \COMMENT{Summarize trajectory}
            \STATE $\text{Summaries}.\text{append}(s_g)$
        \ENDFOR

        \STATE \textcolor{blue}{// Stage 2: Group advantage extraction}
        \STATE $\text{Operations}_i \leftarrow \text{LLM}_{\text{extract}}(\text{Summaries}, \mathcal{E}, q_i, a_i)$
        \STATE $\text{AllOperations}.\text{append}(\text{Operations}_i)$
    \ENDFOR

    \STATE \textcolor{blue}{// Stage 3: Batch consolidation}
    \STATE $\text{FinalOps} \leftarrow \text{LLM}_{\text{consolidate}}(\text{AllOperations}, \mathcal{E})$
    \STATE $\mathcal{E} \leftarrow \text{ApplyOperations}(\mathcal{E}, \text{FinalOps})$

    \STATE \textcolor{blue}{// Evaluation}
    \STATE $\text{Performance}_t \leftarrow \text{Evaluate}(\pi_\theta, \mathcal{E}, \mathcal{Q}_{\text{val}})$
    \STATE \textbf{if} $\text{Performance}_t < \text{Performance}_{t-1}$ \textbf{then} revert $\mathcal{E}$
\ENDFOR
\RETURN $\mathcal{E}$
\end{algorithmic}
\end{algorithm}

\subsection{Stage 1: Trajectory Summarization}

The first stage converts each full trajectory (potentially hundreds of tokens with tool calls, reasoning steps, etc.) into a concise step-by-step summary that highlights:
\begin{itemize}
    \item Actions taken at each step
    \item Which experiences were applied
    \item Where errors or detours occurred (if incorrect)
    \item Core outcomes of each step
\end{itemize}

\textbf{Prompt Template:}
\begin{quote}
\small
\textit{An agent system may be provided with some experiences, and then it produces the following trajectory to solve the given problem. Please summarize the trajectory step-by-step:}

\textit{1. For each step, describe what action is being taken, and which experience has been used in this step.}

\textit{2. Given the grading of this rollout and the correct answer, identify and explain any steps that represent detours, errors, or backtracking.}

\textit{3. Maintain all the core outcome of each step.}

\textit{<trajectory>\{trajectory\}</trajectory>}

\textit{<evaluation>\{correct/wrong\}</evaluation>}

\textit{<groundtruth>\{answer\}</groundtruth>}

\textit{Only return the trajectory summary of each step.}
\end{quote}

\textbf{Example Summary:}
\begin{quote}
\small
\textit{Step 1: Applied experience [G0] to validate geometry solution lies within bounded region. Identified that point $P$ must satisfy $0 \leq x \leq 5$.}

\textit{Step 2: Computed distance formula using Pythagorean theorem. No specific experience applied, standard calculation.}

\textit{Step 3: ERROR - Failed to check discriminant sign before solving quadratic. Should have applied experience [G21] to separate real/imaginary parts first. This led to extraneous solution.}

\textit{Step 4: Backtracked and recomputed with proper validation. Final answer $x = 3$.}
\end{quote}

This summarization serves two purposes:
\begin{enumerate}
    \item \textbf{Compression}: Reduces context length for subsequent LLM calls
    \item \textbf{Attribution}: Explicitly links experiences to outcomes, enabling targeted updates
\end{enumerate}

\subsection{Stage 2: Group Advantage Extraction}

The second stage compares all $G$ trajectory summaries within a group to identify patterns that distinguish successful from unsuccessful attempts. The LLM is prompted to:
\begin{itemize}
    \item Identify key correct decisions in successful trajectories
    \item Pinpoint where and why reasoning went wrong in failed trajectories
    \item Note strategies that were used or missed
    \item Propose up to 3 operations to update the experience library
\end{itemize}

\textbf{Prompt Template:}
\begin{quote}
\small
\textit{Review these problem-solving attempts and extract generalizable experiences:}

\textit{1. Trajectory Analysis:}
\begin{itemize}
    \item[] \textit{- For successful steps: Identify key correct decisions}
    \item[] \textit{- For errors: Pinpoint where/why reasoning went wrong}
    \item[] \textit{- Note patterns or strategies used/missed}
\end{itemize}

\textit{2. Update Existing Experiences:}
\begin{itemize}
    \item[] \textit{- Options: [modify, add, delete]}
    \item[] \textit{- Max 3 operations per group}
    \item[] \textit{- Requirements: Begin with context, focus on strategic patterns}
\end{itemize}

\textit{Return JSON: [\{"option": "add", "experience": "..."\}, ...]}

\textit{<problem>\{problem\}</problem>}

\textit{<trajectories>\{G summaries\}</trajectories>}

\textit{<groundtruth>\{answer\}</groundtruth>}

\textit{<experience>\{current experiences\}</experience>}
\end{quote}

\textbf{Example Extraction:}
\begin{quote}
\small
\begin{verbatim}
[
  {
    "option": "add",
    "experience": "When solving quadratic equations in
     geometry problems, check discriminant sign before
     proceeding to avoid extraneous solutions."
  },
  {
    "option": "modify",
    "experience_id": "G0",
    "new_text": "When solving geometry problems with
     intersections or constrained regions, validate
     solutions satisfy ALL boundary conditions."
  }
]
\end{verbatim}
\end{quote}

\subsection{Stage 3: Batch Consolidation}

The final stage reviews all proposed operations from the entire batch (across all queries) and consolidates them to:
\begin{itemize}
    \item Merge duplicate or highly similar experiences
    \item Ensure each experience is concise ($\leq 32$ words)
    \item Remove low-value or contradictory experiences
    \item Maintain diversity and coverage
\end{itemize}

\textbf{Prompt Template:}
\begin{quote}
\small
\textit{Consolidate suggested updates into final experience revisions:}

\textit{Requirements:}
\begin{enumerate}
    \item[] \textit{1. Clear, generalizable, $\leq$ 32 words}
    \item[] \textit{2. Focus on strategic thinking}
    \item[] \textit{3. Avoid duplication}
\end{enumerate}

\textit{Options: [modify, merge, delete]}

\textit{<experience>\{current\}</experience>}

\textit{<suggested\_updates>\{all group operations\}</suggested\_updates>}

\textit{Return JSON with final operations.}
\end{quote}

\textbf{Example Consolidation:}
\begin{quote}
\small
\begin{verbatim}
[
  {
    "option": "merge",
    "experience_ids": ["temp_1", "temp_5", "temp_12"],
    "new_experience": "For quadratic equations in geometry,
     validate discriminant and boundary conditions before
     accepting solutions.",
    "id": "G45"
  },
  {
    "option": "delete",
    "experience_id": "G23",
    "reason": "Contradicted by empirical evidence in
     current batch."
  }
]
\end{verbatim}
\end{quote}

\subsection{Experience Format and Characteristics}

Each experience follows a consistent format designed for maximum utility:

\begin{itemize}
    \item \textbf{Concise}: $\leq 32$ words to fit many in context window
    \item \textbf{Strategic}: Focus on "when" and "why", not "how to calculate"
    \item \textbf{Contextual}: Begin with triggering condition ("When solving...", "For problems involving...")
    \item \textbf{Actionable}: Specify clear decision or validation step
    \item \textbf{Generalizable}: Avoid problem-specific details, focus on patterns
\end{itemize}

\textbf{Example Experiences from AIME Dataset:}

\begin{enumerate}
    \item[\texttt{[G0]}] \textit{When solving geometry problems with intersections, validate solutions lie within bounded regions or segments, not on extensions, to avoid extraneous answers.}

    \item[\texttt{[G1]}] \textit{For expected extreme statistics in combinatorial problems, use direct enumeration for small sizes.}

    \item[\texttt{[G10]}] \textit{When using mathematical invariants to prove impossibility, always validate them against known achievable states or small cases.}

    \item[\texttt{[G21]}] \textit{For complex polynomials with real parameters, separate real and imaginary parts to find when real roots exist.}

    \item[\texttt{[G37]}] \textit{In geometry problems with points on sides of a triangle and given segment lengths, first determine all three side lengths by summing the appropriate segments.}
\end{enumerate}

\subsection{Context Injection Strategy}

During inference, the experience library is injected into the system prompt:

\begin{quote}
\small
\textit{You are a mathematical problem solver. Use the following learned experiences to guide your reasoning:}

\textit{\# Learned Experiences}

\textit{[G0]. When solving geometry problems with intersections...}

\textit{[G1]. For expected extreme statistics...}

\textit{...}

\textit{Now solve the following problem, explicitly noting which experiences you apply:}

\textit{[Problem statement]}
\end{quote}

For long experience libraries ($>50$ experiences), we use embedding-based retrieval to select only the top-$k$ most relevant experiences for each query:

\begin{equation}
\mathcal{E}_{\text{relevant}}(q) = \text{TopK}\left(\left\{ e \in \mathcal{E} : \text{sim}(\text{embed}(e), \text{embed}(q)) \right\}, k=5\right)
\end{equation}

This balances context richness with inference cost, maintaining the Hamiltonian equilibrium.

